\section{Empirical Design}
\label{sec:empirical-design}

This section sets out the econometric design used throughout the paper. I begin by describing the staggered municipal rollout of biometric voter registration and the panel structure it produces. I then explain why a standard two-way fixed-effects event study is unreliable in this setting and introduce the imputation estimator I use as the baseline. I close by defining the strict and hybrid treatment regimes that delimit the main estimation sample.

\textbf{Staggered rollout.} The empirical design exploits the staggered rollout of biometric voter registration across Brazilian municipalities. Municipalities entered the biometric regime in different election years, beginning with three pilot municipalities in 2008 and accelerating through the 2010, 2012, 2014, 2016, and especially 2018 election cycles. A large set of municipalities remained never treated through the end of the sample. The resulting panel design compares each treated cohort to municipalities that had not entered the biometric regime, with municipality fixed effects absorbing time-invariant local characteristics and election-year fixed effects absorbing national shocks common to all municipalities in a given cycle. The staggered structure is also the reason a more careful estimator is needed. Cohorts that adopt early may differ from cohorts that adopt late not only in pre-treatment levels of the outcome but also in the size, shape, and timing of their treatment response. A specification that cannot accommodate this heterogeneity will conflate it with the underlying dynamic effect.

\textbf{Why not a standard event-study TWFE.} A natural starting point is the dynamic two-way fixed-effects regression
\begin{equation}
\label{eq:twfe-event-study}
Y_{mt} = \alpha_m + \lambda_t + \sum_{k \neq -1} \beta_k \,\mathbf{1}\!\left[t - G_m = k\right] + \varepsilon_{mt},
\end{equation}
where $Y_{mt}$ is the outcome for municipality $m$ in election year $t$, $G_m$ is municipality $m$'s first treatment year, $\alpha_m$ and $\lambda_t$ are municipality and election-year fixed effects, and $k = t - G_m$ indexes event time relative to adoption. I report Equation~\hyperref[eq:twfe-event-study]{(\ref*{eq:twfe-event-study})} for illustration only. Under staggered adoption with heterogeneous treatment effects, the OLS estimator of $\beta_k$ is a weighted average of unit-time treatment effects in which already-treated cohorts serve as implicit controls for later-treated cohorts, and the implicit weights need not be convex and can be negative \citep{GoodmanBacon2018DifferenDifferenWith,deChaisemartin2020TwoWayFixed,Abraham2018EstimatiDynamicTreatmen}. The resulting forbidden comparisons can flip the sign of $\hat\beta_k$ even when every individual treatment effect is positive. The recent staggered-design literature has converged on a class of estimators that avoid these comparisons by restricting controls to not-yet-treated or never-treated units \citep{Callaway2020DifferenDifferenWith,Abraham2018EstimatiDynamicTreatmen,deChaisemartin2020TwoWayFixed,Borusyak2024RevisitiEventStudy}. I use the imputation estimator from this class as the baseline throughout the paper.

\textbf{The imputation estimator.} The baseline estimator is the two-step imputation event-study estimator of \citet{Borusyak2024RevisitiEventStudy}. The estimator posits the same outcome model used implicitly by Equation~\hyperref[eq:twfe-event-study]{(\ref*{eq:twfe-event-study})},
\begin{equation}
\label{eq:bjs-model}
Y_{mt} = \alpha_m + \lambda_t + \tau_{mt}\,D_{mt} + \varepsilon_{mt},
\end{equation}
where $D_{mt}$ equals one for municipality-election cells under biometric treatment and $\tau_{mt}$ is the cell-specific treatment effect. In the first step, I estimate the fixed effects $\hat{\alpha}_m$ and $\hat{\lambda}_t$ by ordinary least squares using only the untreated cells, $\Omega_0 = \{(m,t): D_{mt} = 0\}$. By construction, this first step is uncontaminated by post-treatment outcomes: no treated observation enters the regression that pins down the time path of $\lambda_t$ or the level $\alpha_m$. In the second step, I impute the no-BVR counterfactual for every treated cell as
\[
\widehat{Y}_{mt}(0) = \hat{\alpha}_m + \hat{\lambda}_t,
\]
recover the unit-time treatment effect as $\hat{\tau}_{mt} = Y_{mt} - \widehat{Y}_{mt}(0)$, and aggregate these unit-time effects into an event-time path by averaging $\hat{\tau}_{mt}$ across treated municipalities with the same event time $k = t - G_m$. This procedure solves the negative-weighting problem in Equation~\hyperref[eq:twfe-event-study]{(\ref*{eq:twfe-event-study})}, is robust to heterogeneous treatment effects across cohorts and event times, and is efficient under the homoskedasticity assumption maintained by \citet{Borusyak2024RevisitiEventStudy}.

\textbf{Pre-trends and inference.} The first-step OLS regression also delivers the paper's pre-trend evidence. To obtain pre-period coefficients, I re-estimate the OLS step on $\Omega_0$ with indicators for pre-treatment leads $k \in \{-8, -6, \ldots, -2\}$ included alongside the fixed effects, omitting event time $k = -2$ as the reference period in the plotted paths. The OLS coefficients on those leads are the pre-trend estimates displayed in the event-study figures, and the null hypothesis of no anticipation and parallel pre-trends corresponds to the joint Wald test that those coefficients equal zero \citep{Borusyak2024RevisitiEventStudy}. The pre-period is therefore identified purely from the untreated sample by OLS, while the post-period is identified by imputation from the same first-step model, so the two halves of each event study are estimated on a common set of fixed effects. Standard errors throughout the paper are clustered at the municipality level using the analytic variance estimator in \citet{Borusyak2024RevisitiEventStudy}, which accounts for both the imputation step and the panel structure. Where later sections report Callaway-Sant'Anna or dynamic TWFE event studies alongside the imputation paths, those serve as cross-estimator robustness checks rather than substitutes for the baseline.

\textbf{Strict and hybrid BVR.} The treatment variable in Equation~\hyperref[eq:bjs-model]{(\ref*{eq:bjs-model})} is exposure to the strict biometric regime. Strict BVR refers to municipalities that had fully transitioned to the new system, in which biometric authentication was the operative procedure at the polling place. Hybrid implementation refers to municipalities in which the biometric infrastructure existed and some voters had updated their records, but full enforcement of the strict regime had not yet taken effect. The baseline sample therefore excludes municipalities ever flagged as hybrid, so the treatment contrast is strict BVR against municipalities that had not adopted any biometric regime. This restriction leaves 2{,}480 treated and 1{,}245 never-treated municipalities in the registry-size specification. Appendix~\hyperref[app:including-hybrid]{Section~\ref*{app:including-hybrid}} repeats the main estimates after including hybrid municipalities in the treatment definition. The hybrid-inclusive estimates preserve the sign and shape of the headline results but are attenuated on impact, which is the pattern one would expect if hybrid implementation is a lower-intensity version of the same reform. The empirical sections that follow take this design as given; where additional specification choices are needed, I describe them in place rather than restating the baseline.
